% Provenance: copied from the design workspace CambridgePQ/recursion.fast, `report/sections/07-protocol-model.tex`, on 2026-10-19. Read-only copy; the source is authoritative.
% 07-protocol-model -- written by W2
\section{Protocol model and submission format}\label{sec:protocol}

This section fixes the mathematical objects a submission is about: the proof model, the relation,
the fields, the shapes of protocol rounds, the indexer, the RISC-V machine, and the link between
the machine and the Lean specification. The Lean listings are the statements a reader needs in
context, excerpted from the compiled trusted contract v0.3 (\code{RecursionFast/Contract/}, 796 code
lines; Lean v4.34.1) verbatim except for line breaks, with the file name on the frame;
\enquote{abridged} marks omitted declarations. \Cref{sec:appendix-contract} lists every trusted
file in full. Notation:
$[n]=\{0,\dots,n-1\}$; a \emph{transcript} alternates prover messages and verifier challenges;
``raw vector'' means an element of \lean{Fin (2\^{}logLen) → K}.

\subsection{The proof model: holographic public-coin IOPs}\label{subsec:proto-model}

\begin{decision}{P-1}{Information-theoretic holographic IOP; compilation cited}
Submitters specify a holographic public-coin interactive oracle proof (IOP;
\cref{subsec:bg-iop}) by a Lean \emph{spec verifier}, prove perfect completeness and worst-case
round-by-round knowledge soundness (RBR-KS) of it with no shared oracle, and prove that a RISC-V
image decides exactly as the spec does. Compilation to a succinct non-interactive argument (BCS
with Merkle commitments and Fiat--Shamir, holographic COS variant) is a cited meta-theorem with an
explicit bound (\cref{sec:security}).
\end{decision}

\noindent Three models were considered: (A) an IOP in the information-theoretic model; (B) an
interactive argument in the random-oracle model (ROM), with a hash system call and Merkle paths
checked by the machine, as in sig.golf; (C) a non-interactive Fiat--Shamir argument in the ROM, the
machine being the deployed verifier. (A) is chosen because: (1) ArkLib's sorry-free assets are its
IOP definitions and leanerVM's components, while BCS is absent and Fiat--Shamir soundness is
\lean{sorry} (\cref{subsec:bg-lean}), so under (B) or (C) every submission would redo BCS, estimated
at two to four times sig.golf's 154k lines, against tens of lines of security proof for the
baseline under (A); (2) RBR-KS is what BCS and Fiat--Shamir consume (\cref{thm:bg-rbr-sr,thm:bg-bcs}),
so citing the compiler loses nothing provable today; (3) the batcher is a reduction of knowledge in
the IOP model (\cref{sec:recursion}), whereas in the ROM the inner verifier queries the oracle inside
the proved relation and no theorem applies \cite{BCG24-RelROM}; (4) with no shared oracle every
verifier is a pure function, so the audited text has no probability monad or oracle cache, the
hardest audit items of the golf contracts. Not certified, and listed: Fiat--Shamir attacks on
concrete hashes \cite{KRS25,Fen26-GenR1CS}, hash security, Merkle and transcript implementations,
and recursion after instantiation (\cref{subsec:sec-fs,sec:recursion}).

\paragraph{Security levels.} The roadmap is organized in five security levels L0--L4, distinct from
the score tracks T0--T3 of \cref{sec:cost}, from the library components lib-L0--lib-L6 of
\cref{sec:issues} and from the latency-board stub \code{L1.lean.in}.
\begin{center}\small
\begin{tabularx}{\linewidth}{@{}l>{\raggedright\arraybackslash}p{2.2cm}X@{}}
\toprule
Level & Model & What is new \\
\midrule
L0 & IOP & ``send the witness'' baseline and a RISC-V checker: the inhabitation witness of the certificate \\
L1 & IOP & base track and one-step accumulation; the v1 contract \\
L2 & IOP & tree composition theorem; efficient-extraction track (RISC-V extractor with a proved cycle bound); conditional track with named hypotheses \\
L3 & ROM argument & hash system call, verifier checks Merkle paths; needs a generic, organizer-proved BCS theorem \\
L4 & full recursion & verifier-in-circuit through a trusted RISC-V-to-R1CS compiler; a target schema only \\
\bottomrule
\end{tabularx}
\end{center}

\subsection{The relation}\label{subsec:proto-relation}

The relation is R1CS over the field $K$ of \cref{subsec:proto-fields}, with the constant-one
convention. An \emph{index} is a triple of matrices $(A,B,C)\in (K^{m\times n})^3$, an
\emph{instance} is $x\in K^{n_x}$, a \emph{witness} is $w\in K^{n_w}$ with $n=1+n_x+n_w$, and
\[ \mathsf{R1CS}(i,x,w)\ :\Longleftrightarrow\ (Az)\circ(Bz)=Cz,\qquad z=(1,x,w), \]
with $\circ$ the entrywise product. The constant wire is needed for affine constraints; ArkLib's
R1CS relation has no constant entry and is restated, with an optional bridge for $x_0=1$.

\begin{decision}{P-2}{Fixed class shapes with harness padding}
For $k\in\{16,20,24\}$ the class shape is $m=2^k$ rows and $n=2^k$ columns, $n_x=2^{10}$ public
inputs (so $n_w=2^k-1-2^{10}$), and \emph{sparse} indices: each matrix is a list of entries
$(r,c,v)$, entries at a repeated position are summed, and the three lists have at most
$2^{k+4}$ entries in total (which implies
$|\mathrm{supp}\,A|+|\mathrm{supp}\,B|+|\mathrm{supp}\,C|\le 2^{k+4}$, library lemma
\lean{SpMat.card\_supp\_le}). A smaller R1CS is padded by the harness with the contract's
\lean{SpMat.pad} and \lean{PubInput.pad} (entries moved by a column map, zero public inputs), whose
satisfiability preservation in both directions is a library theorem (\lean{pad\_sat\_iff}).
Certificates are required only at the three class shapes. $N_k:=2^k$.
\end{decision}

\noindent Fixed shapes remove all size casework from succinct proofs (parameters become constants,
errors are discharged by \lean{norm\_num}) and lose nothing, since scores were taken at the caps;
small instances pay for padding, as the cost model already charges. $n_x=2^{10}$ answers A1-M2: at
$2^{16}$, ingesting $x$ would cost 27--50\% of a succinct score identically across architectures; a
benchmark with a large public input must hash it in the circuit. Capping columns at $2^k$ keeps
multilinear designs from padding $z$ to $2^{k+1}$. The relation is non-empty at every class (zero
index, $x=w=0$) and has unsatisfiable instances (the row $1\cdot1=0$ on the constant wire), so
neither completeness nor soundness is vacuous (\cref{subsec:sec-nonvacuity}).

\noindent The index is sparse since v0.3, at the baseline engineer's request: a dense
\lean{Matrix (Fin m) (Fin n) K} at $2^{24}$ has $2^{48}$ entries, so an indexer that must inspect
it cannot run in time linear in $\mathit{nnz}$. The dense reading is a library theorem
(\lean{SpMat.mulVec\_eq}: the sparse product equals Mathlib's \lean{Matrix.mulVec} of the denoted
matrix \lean{SpMat.toMat}; \lean{R1CS\_iff} restates the relation densely). One consequence must be
read correctly: a statement is a \emph{representation}. Two entry lists denoting the same matrices
are different statements with the same truth value, and \lean{VK}, \lean{IDX} may depend on the
representation, so the harness supplies benchmark indices in a fixed canonical form (row-major,
duplicates merged, no explicit zeros), a harness convention outside the contract.

\noindent \code{R1CS.lean} (\cref{sec:appendix-contract}) defines exactly this: \lean{Sizes},
\lean{classSizes}, the sparse \lean{SpMat} with summed duplicates, \lean{Index} with its
\lean{nnz\_le} field, \lean{zVec} $=(1,x,w)$, \lean{R1CS} row by row, and \lean{rel}.

\subsection{The fields \texorpdfstring{$K$}{K} and \texorpdfstring{$E$}{E}}\label{subsec:proto-fields}

\begin{decision}{P-3}{Relation field $K=\mathrm{GF}(2^{64})$, challenge field $E=\mathrm{GF}(2^{192})$}
$K=\mathbb{F}_2[x]/(x^{64}+x^4+x^3+x+1)$, an element encoded in 8 bytes little-endian with bit $i$
the coefficient of $x^i$. $E=K[y]/(y^3+y+1)$, three $K$-limbs (coefficients of $1,y,y^2$), 24
bytes. Both are defined in the contract.
\end{decision}

\noindent The competition is enforced over binary fields (user input U1), and this is leanVM's
production pair \cite{leanVM}, so its circuits, WHIR configuration and leanerVM's Lean components
are reusable. \Cref{subsec:bg-binary} gives the rationale: one $\mathrm{GF}(2^{192})$ challenge per
check is the simplest way to meet the $2^{-128}$ cap, every bit string is a canonical element (so
Fiat--Shamir sampling has no rejection or bias), and additive Reed--Solomon domains exist at every
rate. There are no $K$-valued challenges: a single such challenge has a bad fraction $\ge2^{-64}$, so no
state function meets the cap, and omitting the constructor removes the small-subfield exploit. The
contract's $K$ and $E$ use Lean core only (36 code lines) and $K$ is definitionally equal to CompPoly's
\lean{BF64} \cite{comppoly2026}. \lean{Field} instances, cardinalities, irreducibility (CompPoly's
kernel-checked Rabin certificate for the modulus; an order argument for $y^3+y+1$) and
\lean{rfl}-bridges live in the untrusted library; \lean{K\_spec} (a ring isomorphism to Mathlib's
\lean{AdjoinRoot}) is a control. Zbc is defined as the halves of the same carry-less product, so
RISC-V field routines refine the contract nearly definitionally.

\paragraph{Definitional equality, for readers new to Lean.} \enquote{By \lean{rfl}} and
\enquote{definitionally equal} mean that the kernel reduces both sides to the same normal form; such a
proof trusts nothing beyond the kernel but can be expensive. A \lean{simp} proof instead builds a
chain of rewrites (propositional equality), which a second kernel may check far more slowly
(\cref{sec:baseline}). \lean{decide +kernel} proves a decidable proposition by kernel evaluation
alone; on large literals it costs tens of seconds and gigabytes (\cref{subsec:arch-binding}).

\subsection{Round shapes}\label{subsec:proto-rounds}

\begin{decision}{P-4}{Contract-fixed round shapes; raw message vectors}
A protocol at class $k$ is a list of at most $\mathit{roundCap}=2^{14}$ rounds, each of one of two
contract-fixed shapes: \lean{msg o}, in which the prover sends a raw vector
\lean{Fin (2\^{}o.logLen) → K} (an oracle, read by \code{QUERY} in leaves of $2^{\mathit{logLeaf}}$
elements), and \lean{chal nE idx}, in which the verifier sends $n_E$ uniform elements of $E$ and,
for each $b$ in the list \lean{idx}, a uniform $b$-bit index (\lean{BitVec b}, $b\le 64$ by
admission). Lengths and leaves are powers of two, stored by their exponents. Submitters choose the
shape, not the types. There are no submitter message types, no serializers and no submitter
challenge types.
\end{decision}

\noindent ArkLib, like most of the literature, lets a protocol declare arbitrary message and
challenge types. In a competition this is an attack surface, found independently by the incentive
and the definitions reviews.

\paragraph{The typed-message attack (A1-C1, A2-C1).} Suppose a submitter type \lean{Message j} with
a serializer \lean{ser j : Message j → (Fin len → K)} were allowed. Security is proved on
\emph{typed} transcripts and decision agreement on serialized typed transcripts. A submitter
declares
\[ \text{\lean{Message j}} := \{\,p\in K[X] : \deg p<2^k\,\},\qquad
   \text{\lean{ser j}}\,p := (p(a))_{a\in D}, \]
which is how a polynomial IOP is written in ArkLib, so an agent porting ArkLib components does it
without intending to cheat. In the theorem the prover can only send low-degree polynomials, so the
verifier needs no proximity test: an evaluation claim is checked by a typed quotient and a few
spot checks, about $3.3\times$ cheaper than the best honest design in the incentive model. Every
clause holds, yet the deployed verifier is broken by any non-codeword, on which the machine is
unconstrained. The definitions review proved a starker variant in Lean: the toy ``send the
witness'' protocol with serializer $w\mapsto(1,w)$ and a machine that accepts whenever the flag word
is not 1 satisfies \lean{Secure}, \lean{Agrees} and termination, and accepts the all-zero proof for
an unsatisfiable index.

\paragraph{Why raw vectors fix it.} With \lean{OVal o = Fin (2\^{}o.logLen) → K} and the identity encoding, a typed
transcript \emph{is} an arbitrary assignment of oracle contents, so ``for all transcripts'' is the
quantifier of the IOP model in every theorem we cite \cite[Def.~23.2.1]{CY24-Book},
\cite[Def.~B.1]{BCFW25-WARP}. Structured messages remain possible: the spec verifier parses and
rejects on failure, as a deployed verifier must.

\paragraph{Challenges (A1-M1, A2-C2).} Submitter challenge types fail three ways. Sampling the
submitter type runs candidate code on the trust path. Sampling the codomain of a submitter
serializer produces out-of-image tapes, on which the machine is unconstrained (a \lean{Unit} round
serialized to zeros keeps \lean{Secure} and lets the machine accept on any non-zero word). And any
function of a round's randomness (powers $r^{2^j}$, barycentric weights, domain points, inverses)
becomes free verifier work done by the harness, worth 10--15\% of a succinct score. The contract
menu $E^{n_E}\times\prod_b\{0,\dots,2^b-1\}$ with contract encodings closes all three; the static
Fiat--Shamir surcharge (\cref{sec:cost}) prices the draws.

\paragraph{The round cap (A2-C3).} The interactive error is $\sum_i\eps_i$ and the compiled bound
grows with $k$, so without a round cap the per-round cap bounds nothing. The definitions review
proved in Lean that a protocol with $k$ challenge rounds, no messages and the empty relation, which
accepts iff some challenge is zero, meets the full certificate at the cap for every $k$ (one bad
challenge per round), yet accepts with probability $\ge 1/2$ at $k=|C|$. With $n\le2^{14}$ the
interactive error is at most $2^{-114}$. $2^{14}$ rather than $2^{10}$ admits holographic Spartan
with GKR memory checking (about 1,070 rounds at $2^{24}$); a round costs its Fiat--Shamir surcharge.

\begin{leancode}[fontsize=\footnotesize,label={\footnotesize Contract/Shapes.lean (abridged)}]
/-- Shape of an oracle: `2^logLen` elements of K, read in leaves of `2^logLeaf` elements, so a
QUERY opens one of `2^(logLen - logLeaf)` leaves (Merkle depth `logLen - logLeaf`). -/
structure OShape where
  logLen : ℕ
  logLeaf : ℕ
  deriving DecidableEq

/-- Contents of an oracle of shape `o`. -/
abbrev OVal (o : OShape) := Fin (2 ^ o.logLen) → K

/-- A round of the contract's fixed menu. -/
inductive Round where
  /-- The prover sends an oracle of shape `o`. -/
  | msg (o : OShape)
  /-- The verifier sends `nE` elements of E and one `b`-bit index for each `b ∈ idx`. -/
  | chal (nE : ℕ) (idx : List ℕ)
  deriving DecidableEq

/-- What is sent in a round: a raw K-vector, or (E-vector, indices). -/
def Round.Ty : Round → Type
  | .msg o => OVal o
  | .chal nE idx => (Fin nE → E) × ((j : Fin idx.length) → BitVec (idx.get j))
\end{leancode}
\noindent For a protocol with round shapes \lean{sh : Fin n → Round}, \code{Shapes.lean} also
defines \lean{roundCap} $=2^{14}$, a transcript prefix \lean{Transcript sh m} (the values of rounds
$0,\dots,m-1$), \lean{FullTranscript sh}, the prover and challenge round indices
\lean{MessageIdx sh} and \lean{ChallengeIdx sh}, and \lean{Challenges sh}, a choice of all
challenges (\cref{sec:appendix-contract}). There is no submitter \lean{ProtocolSpec}: a transcript prefix is a dependent function
on \lean{Fin m}, as in ArkLib's \lean{ProtocolSpec.Transcript}, and the untrusted
\code{Bridge/ArkLib.lean} maps a shape to ArkLib's specification definitionally. Every challenge
type is a contract type, finite and non-empty (\code{Library/Challenges.lean}); none is submitted.

\subsection{The indexer}\label{subsec:proto-indexer}

A succinct verifier for an arbitrary R1CS cannot read the matrices; the holographic model
(\cref{subsec:bg-r1cs}) supplies an indexer. A submission provides two \emph{computable} Lean
functions: a key \lean{VK : Index sz → (Fin vkLen → K)}, $\mathit{vkLen}\le 2^{10}$, placed on the
tape, and an index oracle \lean{IDX : Index sz → OVal idxShape}, oracle 0, reachable only by
priced queries. The spec verifier never sees $i$, and security is stated for the derived verifier
$((i,x),\mathrm{tr})\mapsto\mathsf{verify}(\mathsf{VK}\,i,x,\mathsf{IDX}\,i,\mathrm{tr})$ over
\emph{every} index of class shape, adversarial ones included. After compilation the key is extended
by the Merkle root of \lean{IDX i}. The \emph{key cap} is load-bearing: the key is readable at
random-access cost and is an arbitrary function of $i$, so uncapped, \lean{VK i := (A,B,C)} defeats
holography and is a free lookup table. The index oracle needs only the admission length limit
$\mathit{idxLen}\le2^{k+8}$, since every access is priced; the limit also bounds the side condition
of the compiled bound (\cref{subsec:sec-cap}) [A5-M4]. \emph{Feasibility} of the indexer (A2-M4)
is not a Lean property. \lean{IDX} could be the table $x\mapsto[\exists w,\ \mathsf{R1CS}(i,x,w)]$,
which is computable in Lean's sense (enumeration over the finite $K^{n_w}$) but infeasible, and
would make the certified key undeployable. Lean's \lean{noncomputable} marks definitions without
compiled code (choice, real inverses), not expensive ones. The environment gate therefore only
guarantees that \lean{VK}, \lean{IDX}, the prover and the extractor \emph{have} compiled code, and
rejects any \lean{implemented\_by}, \lean{extern}, \lean{csimp}, \lean{init}, \lean{builtin\_init}
or \lean{export} entry registered by an untrusted module (\cref{subsec:arch-ci}); other fields,
such as the errors \lean{ε}, may be \lean{noncomputable}. Feasibility is enforced by the badge path, which runs the compiled indexer
on hidden indices under a time and memory budget and gates record eligibility, not the proved score
(\cref{subsec:arch-badge}). In the
contract the index-oracle cap is not a structure field but an admission clause of the certificate
(\lean{OShape.IdxAdmissible}: $2^{\mathit{logLen}}\le 2^8\cdot\max(m,\mathit{cols})$, i.e.\
$\mathit{idxLen}\le2^{k+8}$ at class $k$; message oracles keep $2^{k+10}$).

\begin{leancode}[fontsize=\footnotesize,label={\footnotesize Contract/Protocol.lean (abridged)}]
/-- A holographic public-coin IOP for `rel sz`. -/
structure Protocol (sz : Sizes) where
  /-- Number of rounds, at most `roundCap = 2^14`. -/
  n : ℕ
  n_le : n ≤ roundCap
  /-- The round shapes (contract menu). -/
  shape : Fin n → Round
  /-- Verification key length and the (honest, computable) key. -/
  vkLen : ℕ
  vkLen_le : vkLen ≤ vkCap
  VK : Index sz → (Fin vkLen → K)
  /-- The index oracle (oracle 0): its shape and the (honest, computable) indexer. -/
  idxShape : OShape
  IDX : Index sz → OVal idxShape
  /-- The honest prover. -/
  prover : Strategy shape (Index sz × PubInput sz) (Wit sz) Unit Unit
  /-- The spec verifier's decision, from `VK i`, `x`, `IDX i` and the full transcript. -/
  verify : (Fin vkLen → K) → PubInput sz → OVal idxShape → FullTranscript shape → Bool
  /-- The named round-by-round extractor. -/
  WitMid : Fin (n + 1) → Type
  extractor : Extractor shape (Index sz × PubInput sz) (Wit sz) Unit WitMid
  /-- Declared per-round errors. -/
  ε : ChallengeIdx shape → ℝ≥0
\end{leancode}
\noindent \lean{vkCap} is $2^{10}$; \lean{specVerifier} accepts iff \lean{verify (VK i) x (IDX i) tr},
and \lean{Protocol.Secure} is given in \cref{subsec:sec-rbrks}.


\subsection{The machine}\label{subsec:proto-machine}

\begin{decision}{P-5}{One contract-owned oracle machine}
RV64IM with all W-forms, decoded by the contract from raw 32-bit words, plus Zbc (\code{clmul},
\code{clmulh}, \code{clmulr}). Harvard code space at \lean{CODE\_BASE}${}=\mathtt{0x1000}$; flat
data memory of $2^{36}$ bytes (a function in Lean, sparse in the emulator); natural alignment;
every fault is a rejection. System calls (\code{ECALL}) are selected by register \code{t0}:
\begin{itemize}[nosep]
\item \code{0 HALT}: accept iff $a_0=1$ and no child hart is live; $a_0=0$ rejects; any other value
  traps.
\item \code{1 QUERY}: $a_0=j$ selects an oracle, $a_1$ a leaf index, $a_2$ an 8-aligned writable
  destination. The call writes leaf $a_1$ of oracle $j$ ($2^{\mathit{logLeaf}}$ elements of $K$,
  one doubleword each) at $a_2$; an out-of-range oracle, leaf or destination traps.
\item \code{2 HASH}: reserved for the ROM track; it traps in v0.3, as does every unassigned value.
\item \code{3 SPAWN}, \code{4 JOIN}: fork-join harts, only in an image that sets its \lean{harts}
  flag; otherwise they trap exactly as in v0.2.
\end{itemize}
Memory map: image data words at address~0; the read-only tape window
$[2^{32},2^{32}+2^{20})$; the T2 output accumulator at \lean{OUT\_BASE}${}=2^{33}$;
$\mathit{sp}=2^{36}$. The tape holds the four sizes $(m,n_x,n_w,\mathit{nnz})$ $\|$ \lean{VK i}
$\|$ $x$ $\|$ the challenges in round order, one word per $K$-element ($E$ as three limbs, an index
zero-extended). The oracles are \lean{IDX i} ($j=0$) and the prover messages ($j\ge1$), reachable
only through \code{QUERY}. Each hart runs for at most $F=2^{36}$ steps, and \textbf{fuel exhaustion
is a rejection}.
\end{decision}

\noindent There is exactly one machine: during design three incompatible ones existed (baseline
assembly, design text, toolchain probe), and every measurement against them had to be redone
(A3-C1). The base instruction semantics and decoder are riscv-zkvm's hand model, vendored
byte-identically (\lean{execInstrBr}, \lean{decode}); the contract adds the W-op and Zbc decoder arms
and semantics, memory checks, system calls, loading, fuel, harts and accounting. The choices are
justified as follows.
\begin{itemize}[nosep]
\item \emph{Zbc.} A $K$ product costs about 379 cycles with table methods and 35 with Zbc; a 1\,MiB
  image cannot buy this back. Zbc has its own counter, priced per zkVM profile, so re-weighting
  needs no new cost proof (\cref{sec:cost}).
\item \emph{$2^{36}$ bytes.} At the $2^{24}$ caps the baseline's index alone is about 4.4\,GiB. The
  Lean memory is a function, so the bound is free in proofs; the emulator allocates lazily.
\item \emph{Leaf-returning QUERY, closed form.} Single-element queries would make a $2^4$-fold query
  pay sixteen Merkle paths, and an amortized per-element price is gamed by huge leaves read once.
  A declared leaf priced as one path plus its bytes (\lean{CostVec.addQuery}: one query,
  $\mathit{logLen}-\mathit{logLeaf}$ levels, $8\cdot2^{\mathit{logLeaf}}$ bytes) is exact for honest
  designs. The memory effect is a closed form (the 8 bytes at $a_2+8t$ hold leaf element $t$; all
  else unchanged), not a fold of stores, so readback lemmas need no induction over the leaf (A3-C3).
\item \emph{Challenges on the tape.} The image is a plain program on a static input, and the
  theorems quantify over challenges like any other input.
\item \emph{Fuel exhaustion rejects; no termination theorem} (A3 §6.1). \lean{execute} returns
  within $F$ steps per hart by construction. This bounds denial of service at $2^{36}$ steps for an
  image without harts and at $(\#\mathrm{spawns}+1)\cdot2^{36}$ with harts, since a child runs on its
  parent's remaining fuel; a deployment must copy this rule exactly (A6-F1, v0.4 gap). \lean{Agrees} forces every spec-accepted (in particular every honest) transcript to be accepted
  within fuel. A termination theorem over all raw tapes would need a ranking argument per
  data-dependent loop, at no gain for recursion. The deployed semantics includes the cap.
\item \emph{Decision separate from accounting.} \lean{execute} threads the accounting \lean{Acct}
  (the 18 counters of \cref{subsec:cost-vector}, with first-touch sets of doublewords and 1\,KiB
  pages, code and data keyed apart, a written page counted twice) and the scoreboard timer
  \lean{SB} next to the machine state. No branch of the decision reads them, except the spawn
  counter that issues handles and enforces the $2^{16}$ cap. For an image without harts the
  decision and final state are v0.2's word for word and the counters coarsen to v0.2's
  (\code{Library/V02.lean}), so v0.2 \lean{Agrees} proofs transfer; \code{Library/Seq.lean} gives a
  timer-free single-hart machine equal to \lean{run}, so cost proofs never mention the timer.
\item \emph{Fork-join harts, sequentialised.} \code{SPAWN} ($a_0$ entry, $a_1$ argument, $a_2$
  child stack pointer, $a_3/a_4$ write region $W_r$, $a_5/a_6$ read region $R$) runs the child to
  completion at once, on the parent's remaining fuel. The child loads only from $R\cup W_r$ and
  stores and queries only into $W_r$ (its stack lies in $W_r$); its writes reach the parent at
  \code{JOIN}. While a child is live, no other hart stores into its $R\cup W_r$ or loads from its
  $W_r$. Every violation, every failing child and \code{HALT} with a live child rejects, so every
  interleaving reaches the same state at \code{JOIN} and a deployed runtime may run children in
  parallel (argued, not proved in Lean; A6-N5). Caps: depth 8, 64 live children per hart, $2^{16}$ spawns. Work and span are defined in
  \cref{subsec:cost-lat}.
\end{itemize}

\noindent \code{Machine.lean} (\cref{sec:appendix-contract}) holds the constants above, the
\lean{Image} (code words, data words, \lean{harts} flag), the 18-field \lean{CostVec} with the
\code{QUERY} charge \lean{CostVec.addQuery} (one query, $\mathit{logLen}-\mathit{logLeaf}$ levels,
$8\cdot2^{\mathit{logLeaf}}$ bytes), the decoder, \lean{execute} (one hart, threading the timer
\lean{SB} and the accounting \lean{Acct}), and the two entry points used below: \lean{runFull},
which returns the counters, the pipeline work and the span of an accepting run, and \lean{run},
which returns the counters and the final state.

\subsection{The verifier--spec link}\label{subsec:proto-agrees}

The link between the image and the Lean spec is one equation, quantified as follows.

\begin{definition}[Decision agreement]\label{def:proto-agrees}
An image $\mathit{img}$ \emph{agrees} with a protocol $P$ at class $k$ if for every index $i$ of
class shape $k$, every public input $x$, every raw oracle content of every \lean{msg} round and
every challenge vector, with $\mathrm{tr}=(\vec\Pi,\rho)$,
\[ \mathsf{isSome}\,\mathsf{run}\bigl(\mathit{img},\ \mathsf{sizes}\,\|\,\mathsf{VK}\,i\,\|\,x\,\|\,\rho,\
   [\mathsf{IDX}\,i,\Pi_1,\Pi_2,\dots]\bigr)
   \;=\; \mathsf{verify}(\mathsf{VK}\,i,\ x,\ \mathsf{IDX}\,i,\ \mathrm{tr}). \]
\end{definition}

\noindent The preprocessing is \emph{honest}: tape key and oracle 0 are the submission's own
indexer outputs on an index of the class, computed by the harness, whose root enters the
Fiat--Shamir prefix (\cref{subsec:sec-fs}). Everything the prover or the environment controls is
universally quantified. Two earlier formulations failed in opposite directions. Quantifying over
typed transcripts was too weak (the typed-message attack). Quantifying over arbitrary raw tapes and
keys was too strong: loop counts come from the key and the index, so with adversarial index data a
loop walks to the end of memory, the baseline's cost bound becomes false, and every loop needs a
ranking argument (A3-C2). One equation gives both of ots.golf's directions, \emph{soundness of the
image} (machine accepts $\Rightarrow$ spec accepts, excluding accept-all) and \emph{faithfulness}
(the converse, excluding reject-all).

\paragraph{Costs are never transferred.} \lean{Agrees} is weaker than equality of query trees: the
image may query in any order and number and memoize. This is sound because all costs are read off
the machine (\cref{sec:cost}), and necessary because decision agreement cannot transfer cost: a
spec that queries $H(x)$ twice and a machine that queries once decide identically under a cache
\cite{otsgolf2026}.

\begin{leancode}[fontsize=\footnotesize,label={\footnotesize Contract/Agrees.lean (abridged)}]
namespace Protocol
variable {sz : Sizes} (P : Protocol sz)

/-- The tape: `sizes ‖ VK i ‖ x ‖ challenges` (challenges in round order, `Round.tapeWords`). -/
def tape (i : Index sz) (x : PubInput sz) (tr : FullTranscript P.shape) : List Word :=
  sz.header ++ wordsOf (P.VK i) ++ wordsOf x ++ tr.tapeWords

/-- The QUERY oracles: `IDX i` is oracle 0, the prover messages are oracles `1, 2, …`. -/
def oracles (i : Index sz) (tr : FullTranscript P.shape) : List Oracle :=
  ⟨P.idxShape, P.IDX i⟩ :: tr.oracles

/-- **Decision agreement** of the image with the spec verifier, on all transcripts. -/
def Agrees (img : Image) : Prop :=
  ∀ (i : Index sz) (x : PubInput sz) (tr : FullTranscript P.shape),
    (run img (P.tape i x tr) (P.oracles i tr)).isSome = P.verify (P.VK i) x (P.IDX i) tr

/-- The deployed machine as a pure verifier on `(i, x)` (for the transfer lemma,
`Library/Transfer.lean`: `Agrees` + `Secure` ⇒ the machine verifier is secure). -/
def machineVerifier (img : Image) : PureVerifier P.shape (Index sz × PubInput sz) Unit :=
  fun s tr => (run img (P.tape s.1 s.2 tr) (P.oracles s.1 tr)).map fun _ => ()

end Protocol
\end{leancode}
\begin{leancode}[fontsize=\footnotesize,label={\footnotesize Contract/Cost.lean (abridged)}]
/-- Componentwise order on counters. -/
instance : LE CostVec := ⟨fun a b => List.Forall₂ (· ≤ ·) a.toList b.toList⟩

-- …

/-- **Accepting-run cost bound**: every accepting run, for every index of the size, every `x`
and every transcript, has counters `≤ c` componentwise. -/
def Protocol.AcceptingCostAtMost {sz : Sizes} (P : Protocol sz) (img : Image) (c : CostVec) : Prop
    :=
  ∀ (i : Index sz) x (tr : FullTranscript P.shape) v s,
      run img (P.tape i x tr) (P.oracles i tr) = some (v, s) → v ≤ c

/-- **Accepting-run latency bound** (latency board,
    optional): pipeline work `≤ W` and span `≤ S`. -/
def Protocol.AcceptingLatencyAtMost {sz : Sizes} (P : Protocol sz) (img : Image) (W S : ℕ) : Prop
    :=
  ∀ (i : Index sz) x (tr : FullTranscript P.shape) f s,
    runFull img (P.tape i x tr) (P.oracles i tr) = some (f, s) → f.work ≤ W ∧ f.span ≤ S
\end{leancode}


\noindent Because \lean{tr} ranges over raw vectors and contract challenge types,
``$\forall$ \lean{tr}'' is ``for all oracle contents and all challenges'', and the one certified
vector $c_k$ bounds every accepting run of the class.

\paragraph{Security transfer.} The verifier enters the security definitions only through its
decision, in the completeness equation and in the \textsf{full} law of the knowledge state function
(\cref{sec:security}). Hence:

\begin{lemma}[Transfer along \lean{Agrees}]\label[lemma]{lem:proto-transfer}
If $\mathit{img}$ agrees with $P$ and $P$ is secure, then the machine verifier
$(i,x),\mathrm{tr}\mapsto\mathsf{isSome}\,\mathsf{run}(\dots)$ (\lean{Protocol.\allowbreak machineVerifier})
is perfectly complete and RBR-KS with the same errors, extractor and state function
(\lean{Protocol.\allowbreak secure\_machine}, untrusted library \code{Library/\allowbreak Transfer.lean}; the same for
reductions, \lean{Reduction.secure\_machine}).
\end{lemma}

\noindent The proof is a few lines: by \lean{Agrees} and function extensionality the machine
verifier \emph{equals} the spec verifier. What makes the lemma meaningful is the raw-vector
quantifier: the object BCS compiles is the machine on arbitrary strings, and with contract-fixed
shapes that is exactly the object the lemma is about.

\subsection{T0: the query-program midpoint}\label{subsec:proto-qprog}

The first season is Lean-only (score track T0), so that protocol and security work can start
before RISC-V proof tooling exists. A pure spec verifier has no notion of \enquote{number and price
of queries}; ArkLib's query count is \lean{sorry}, and VCVio's query bounds count calls rather than
depth-weighted prices and bring the probability monad back. The contract therefore adds
\lean{QProg}, a free monad over \code{QUERY} with free local computation (21 code lines).

\noindent \lean{QProg} has two constructors, \lean{done a} and \lean{query j l k} (read leaf $l$ of
oracle $j$ and continue with it); \lean{QProg.eval} runs it against the oracles and
\lean{QProg.cost} adds the query counters along the executed path (listing in
\cref{sec:appendix-contract}). Queries may be adaptive, and \lean{QProg.\allowbreak cost} adds exactly what the machine's
\code{QUERY} adds (\lean{CostVec.\allowbreak addQuery}). A T0 claim therefore has only the three query
counters (all others are $0$: local computation is free), and the T0 score is the recursion-board
formula on them, a lower bound for the T1 score of an image that issues at least the query
program's queries (\lean{Agrees} does not transfer query counts, so an image that memoises may
score lower). The same \lean{QProg} is the midpoint of the T1 refinement. The protocol-level clause
(\lean{agrees} of \lean{Certificate0}) is proved once, with no RISC-V; the machine-level proof then
shows that the image equals the query program between system calls, which is where symbolic
execution stops anyway (no library lemma does this yet). T0 protocol and security proofs therefore
carry over into T1.

\subsection{The certificate}\label{subsec:proto-certificate}

A base-track (T1) submission is a protocol per class (\lean{Submission1}). The image is \emph{not}
part of the submission: like the declared shape summaries, the claimed counters and the claimed
score, it is organiser data rendered into the challenge from the submitted files
(\lean{Declared1}; A4-F2, \cref{subsec:arch-binding}). For each class $c$ (exponent
$k\in\{16,20,24\}$), \lean{Certificate1} requires: \lean{summary}, the protocol's shape summary
(rounds, message rounds, $E$-challenges, index challenges, total message length, tape input words)
equals the rendered one; \lean{score}, the contract's \lean{scoreZk4} of the rendered summaries and
claims equals the rendered \lean{zk4}; \lean{secure} (\cref{sec:security}); \lean{admissible}, the
decidable caps (message oracles $\le2^{k+10}$ elements with leaf $\le$ length, index challenges
$\le64$ bits, tape $\le2^{20}$ bytes, blow-up $V_k\le 2^8$, index oracle $\le2^{k+8}$, at most
$2^{16}$ queries; $n\le2^{14}$ and $\mathit{vkLen}\le2^{10}$ are fields of \lean{Protocol});
\lean{imageSize}, code plus data $\le2^{20}$ bytes; \lean{agrees}; and \lean{cost}, the accepting-run
bound by the rendered claim. The stub states one theorem, \lean{certified}, in which
\lean{submission} is the only definition the solution provides (comparator's
\code{definition\_names}). Comparator thus compares one statement whose every scalar is organiser
data, and the scorer evaluates nothing from the solution: it takes the 48th root of the certified
\lean{zk4}. T0 (\lean{Certificate0}) replaces the image by a query program per class; T2
(\lean{Certificate2}, \cref{sec:recursion}) has three images and a rendered table of
$S^4$ triples. \lean{Certificate2} has no decider cap (DESIGN §15.3, A5-M8, superseding the
earlier $C_{\mathrm{dec}}\le2^{32}$): the decider is bounded only by the fuel and enters the T2
score at weight $1/64$, so linear (Arc/WARP-style) deciders are admissible. It has no all-runs
clause either; an all-runs bound for deciders is planned, not part of v0.3. The latency board adds
the optional \lean{Latency1} on top of \lean{Certificate1}; the baseline's latency record is the
proof of this clause.

\begin{leancode}[fontsize=\footnotesize,label={\footnotesize Contract/Certificate.lean (abridged)}]
/-- A T1 submission: a protocol per class (the image is rendered challenge data). -/
structure Submission1 where
  protocol : (c : Fin 3) → Protocol (classSz c)

/-- Rendered T1 data: the image, shape summaries, claimed counters, claimed `scoreZk4`. -/
structure Declared1 where
  image : Image
  summary : Fin 3 → Summary
  claim : Fin 3 → CostVec
  zk4 : ℚ

/-- T1 certificate (DESIGN §4, §5, §15). -/
structure Certificate1 (S : Submission1) (D : Declared1) : Prop where
  summary : ∀ c, (S.protocol c).summary = D.summary c
  score : scoreZk4 D.summary D.claim = D.zk4
  secure : ∀ c, (S.protocol c).Secure
  admissible : ∀ c, (S.protocol c).Admissible (D.claim c)
  imageSize : D.image.byteSize ≤ imageCap
  agrees : ∀ c, (S.protocol c).Agrees D.image
  cost : ∀ c, (S.protocol c).AcceptingCostAtMost D.image (D.claim c)

/-- Latency-board record (on top of `Certificate1`): claimed `(W, S)` per class, `scoreLat3`. -/
structure Latency1 (S : Submission1) (D : Declared1) (lat : Fin 3 → ℕ × ℕ) (lat3 : ℚ) : Prop where
  bound : ∀ c, (S.protocol c).AcceptingLatencyAtMost D.image (lat c).1 (lat c).2
  score : scoreLat3 lat = lat3
\end{leancode}

\begin{leancode}[fontsize=\footnotesize,label={\footnotesize Challenge/T1.lean.in (abridged)}]
{{image:image}}                  -- rendered: code/data chunks decoded by the contract's unpack
def summary : Fin 3 → Summary := ![⟨{{nat:c16.n}}, {{nat:c16.msgs}}, …⟩, …]
def claim : Fin 3 → CostVec := ![{ add := {{nat:c16.add}}, …, qBytes := {{nat:c16.qBytes}} }, …]
/-- The rendered record; `zk4 = Score_zk⁴⁸`. -/
def declared : Declared1 := ⟨image, summary, claim, {{rat:zk4}}⟩
/-- The submission: a protocol per class (provided by the solution). -/
def submission : Submission1 := sorry
theorem certified : Certificate1 submission declared := sorry
\end{leancode}
